\documentclass[10pt]{article}
\usepackage{times}
\usepackage[margin=2.5cm]{geometry}
\usepackage{booktabs}

\title{User Manual}
\author{}
\date{}

\begin{document}
\maketitle

\section{Overview}

This tool compares a genetic algorithm (GA) against random search for discovering
fairness violations in machine learning classifiers.  All interaction is through
command-line scripts.

\section{Directory structure}

\begin{verbatim}
fairness-ga/
  src/                    Core library modules
    config.py             Default hyperparameters and dataset configs
    preprocessing.py      Data loading and encoding
    model.py              Classifier training
    fairness.py           Discrimination score function
    genetic_algorithm.py  GA implementation
    random_search.py      Random search baseline
    metrics.py            Fairness metric functions
    utils.py              Feature-range sampling utilities
  experiments/
    run_experiment.py     Main experiment driver
    statistical_test.py   Wilcoxon signed-rank tests
    visualization.py      Result plotting
    results/              Generated output (JSON, NPZ, PNG)
  scripts/
    download_adult.py     Download Adult dataset
    download_benchmark_datasets.py
                          Download COMPAS and German Credit
  data/                   Downloaded CSV files
\end{verbatim}

\section{Running experiments}

\subsection{Basic usage}

\begin{verbatim}
python experiments/run_experiment.py
\end{verbatim}

This runs all three datasets with default settings (30~trials, budget~50\,000).

\subsection{Command-line arguments}

\begin{table}[h]
\centering
\small
\begin{tabular}{lll}
\toprule
Argument & Default & Description \\
\midrule
\texttt{--trials}     & 30    & Number of independent trials \\
\texttt{--budget}     & 50000 & Model evaluations per method per trial \\
\texttt{--datasets}   & adult,compas,german & Comma-separated dataset keys \\
\texttt{--ga-population} & 50 & GA population size \\
\texttt{--ga-mutation}   & 0.1 & GA mutation rate \\
\texttt{--with-sensitivity} & off & Run hyperparameter sensitivity sweep \\
\texttt{--sensitivity-trials} & 5 & Trials per sensitivity configuration \\
\texttt{--sensitivity-pop-sizes} & 10,20,40 & Population sizes to test \\
\texttt{--sensitivity-mutation-rates} & 0.05,0.1,0.2 & Mutation rates to test \\
\bottomrule
\end{tabular}
\end{table}

\subsection{Example: quick test}

\begin{verbatim}
python experiments/run_experiment.py \
    --datasets adult --trials 3 --budget 5000
\end{verbatim}

\section{Statistical testing}

After the experiment finishes:

\begin{verbatim}
python experiments/statistical_test.py
\end{verbatim}

Reads \texttt{experiments/results/combined\_results.json} and prints the Wilcoxon
signed-rank test results to the console and to
\texttt{experiments/stat\_test\_summary.txt}.

\section{Visualization}

\begin{verbatim}
python experiments/visualization.py --dataset adult
python experiments/visualization.py --dataset compas
python experiments/visualization.py --dataset german
\end{verbatim}

Each command generates a four-panel figure saved as
\texttt{experiments/fairness\_ga\_results\_<dataset>.png}.

\section{Output files}

\begin{table}[h]
\centering
\small
\begin{tabular}{ll}
\toprule
Path & Contents \\
\midrule
\texttt{results/combined\_results.json} & All metrics (JSON) \\
\texttt{results/<dataset>/trial\_metrics.npz} & Per-trial metric arrays \\
\texttt{results/<dataset>/convergence.npz} & Per-generation convergence curves \\
\texttt{results/<dataset>/sensitivity.json} & Sensitivity sweep results \\
\texttt{stat\_test\_summary.txt} & Wilcoxon test summary \\
\texttt{fairness\_ga\_results\_<dataset>.png} & Visualization figures \\
\bottomrule
\end{tabular}
\end{table}

\end{document}
